\chapter{Kesimpulan dan Saran}

\section{Kesimpulan}

Berisi kesimpulan dari hasil penelitian yang telah dilakukan. Kesimpulan harus menjawab rumusan masalah dan tujuan penelitian yang telah ditetapkan pada Bab 1.

Kesimpulan ditulis dalam bentuk poin-poin atau paragraf yang ringkas dan jelas. Hindari menambahkan informasi baru yang tidak dibahas pada bab-bab sebelumnya.

Contoh format kesimpulan:
\begin{enumerate}
    \item Kesimpulan pertama yang menjawab rumusan masalah pertama.
    \item Kesimpulan kedua yang menjawab rumusan masalah kedua.
    \item Dan seterusnya.
\end{enumerate}

\section{Saran}

Berisi saran untuk pengembangan penelitian selanjutnya atau perbaikan yang dapat dilakukan. Saran dapat berupa:
\begin{itemize}
    \item Saran untuk penelitian lanjutan
    \item Saran untuk perbaikan metode atau sistem
    \item Saran untuk implementasi praktis
    \item Keterbatasan penelitian yang perlu diatasi di penelitian mendatang
\end{itemize}

\textcolor{red}{Catatan: Pastikan kesimpulan dan saran sesuai dengan hasil penelitian yang telah dilakukan dan dibahas pada bab-bab sebelumnya.}
